\documentclass[11pt]{article}

\usepackage[a4paper,margin=27mm]{geometry}
\usepackage[T1]{fontenc}
\usepackage{lmodern}
\usepackage{microtype}
\usepackage{amsmath,amssymb,mathtools}
\usepackage{array,booktabs,longtable,tabularx}
\usepackage{enumitem}
\usepackage{xurl}
\usepackage[hidelinks]{hyperref}

\setlength{\emergencystretch}{3em}
\newcolumntype{Y}{>{\raggedright\arraybackslash}X}
\newcolumntype{L}[1]{>{\raggedright\arraybackslash}p{#1}}

\newcommand{\Ea}{E_{a}}
\newcommand{\BZ}{\mathcal{B}}
\newcommand{\ii}{\mathrm{i}}
\newcommand{\Rc}{R_{\mathrm{c}}}
\newcommand{\Hret}{H_{\mathrm{ret}}}

\title{Research Protocol:\\Periodic-1D Rank-One Bloch-Hopping Models}
\author{Project Koios}
\date{Calculation-owned protocol, version 1}

\begin{document}
\maketitle

\begin{abstract}
This protocol defines the construction, evaluation, and independent verification of
periodic one-dimensional rank-one Bloch-hopping models.  The primary Fourier duality is
the ordinary-lattice displacement--crystal-momentum pair $R$--$k$.  The protocol binds
physical lattice geometry, physical energy units, the lattice kinetic-energy scale
$\Ea=\hbar^2/(m_ea^2)$ without a factor $1/2$, exact Hermiticity, explicit conversion
between finite cyclic coefficients and ordinary-lattice hoppings, and separate
interpolation, support-truncation, and total error channels.  Rank one denotes retained
state-space dimension one, not algebraic matrix rank.  The protocol defines a finite
numerical control and reduced-model contract; it does not establish continuum
convergence, material validity, fitting optimality, or publication-qualified evidence.
\end{abstract}

\noindent\textbf{Keywords:} Bloch Hamiltonian; hopping model; Fourier duality;
rank-one retained space; support truncation; independent verification

\tableofcontents
\clearpage

\section{Protocol status}

\begin{center}
\small
\begin{tabularx}{\textwidth}{L{0.34\textwidth}Y}
\toprule
\textbf{Field} & \textbf{Value} \\
\midrule
Protocol identity & \texttt{periodic1d-rank1-rk-bloch-hopping-v1} \\
Calculation owner & \path{calculations/periodic1d_rank1/} \\
Model protocol owner & \path{calculations/periodic1d_rank1/model/protocol/} \\
Scientific dimension & One periodic lattice dimension \\
Retained-space dimension & One \\
Primary Fourier pair & Ordinary-lattice displacement $R$ and Bloch momentum $k$ \\
Scientific model role & \texttt{REDUCED\_MODEL} \\
Current retained calculation & One predecessor local unsealed finite control \\
New protected or external execution authorized & No \\
Continuum, material, validation, or publication claim & None \\
\bottomrule
\end{tabularx}
\end{center}

This document is the model-level protocol for the operation-first rank-one context.  It
does not modify the bytes, identities, or interpretation of the retained predecessor
execution at
\path{calculations/periodic1d_rank1/band_equivalence_hopping_locality/executions/direct-controlled-four-state-n256/}.
A future execution governed by this protocol must use a new identity and must retain its
own configuration, result, verification, and visualization digests.

\section{Research objective}

The study asks how a scalar periodic Bloch Hamiltonian is represented by
ordinary-lattice hopping coefficients, how that representation depends on finite Fourier
construction, and what changes when hopping support is deliberately restricted.  The
research questions are:

\begin{enumerate}[leftmargin=*]
  \item Does the declared complete finite Fourier representation reconstruct the source
    scalar field on its construction mesh?
  \item How accurately does that complete representation interpolate the source field on
    a disjoint evaluation mesh?
  \item How is complete hopping magnitude distributed over ordinary-lattice
    displacements?
  \item What support is removed at each declared finite range, and what error is caused
    specifically by that removal?
  \item Does the total finite-range candidate satisfy the declared application threshold
    and mesh-stability rules?
  \item Which alignment questions are inapplicable because the represented operator is
    scalar and the internal gauge group is $U(1)$?
\end{enumerate}

The study keeps representation conversion, interpolation, descriptive locality,
support removal, refitting when present, total error, and feature diagnostics separate.
No one channel is substituted for another.

\section{Scientific objects and notation}

\subsection{Direct and reciprocal spaces}

Let $a>0$ be the one-dimensional lattice constant.  The ordinary direct lattice is
\begin{equation}
  \Lambda=\{R=na:n\in\mathbb Z\}.
\end{equation}
The first Brillouin zone is represented by
\begin{equation}
  \BZ=[-\pi/a,\pi/a).
\end{equation}
The primary Fourier pair is
\begin{align}
  H(k) &= \sum_{R\in\Lambda} T(R)e^{\ii kR},
  \label{eq:forward-bloch}\\
  T(R) &= \frac{a}{2\pi}\int_{\BZ}H(k)e^{-\ii kR}\,\mathrm{d}k.
  \label{eq:inverse-hopping}
\end{align}
Here $R$ is an ordinary-lattice displacement, $k$ is Bloch crystal momentum, $H(k)$ is
the scalar represented Hamiltonian, and $T(R)$ is its hopping coefficient.

The dimensionless reduced coordinate $q$ may be used only when explicitly defined by
\begin{equation}
  q=\frac{ak}{2\pi}\in[-1/2,1/2).
\end{equation}
For $R=na$, $e^{\ii kR}=e^{2\pi\ii qn}$.  The reduced coordinate is a derived
coordinate convention; it does not replace the physical $R$--$k$ statement of the
model.

\subsection{Meaning of rank one}

Rank one means that the retained state space at every $k$ has dimension one.  It does
not mean that a sampled matrix has algebraic rank one, that the ambient parent has one
state, or that only one hopping displacement is retained.  In the retained basis,
$H(k)$ and every $T(R)$ are scalar complex quantities.

\subsection{Hermiticity}

The Bloch Hamiltonian is Hermitian for every $k$ exactly when
\begin{equation}
  T(-R)=T(R)^{*}
  \label{eq:scalar-hermiticity}
\end{equation}
for every displacement and $T(0)\in\mathbb R$.  A finite-range constructor therefore
accepts the onsite coefficient and positive-displacement coefficients, then generates
the negative-displacement conjugate partners.  Callers must not independently supply
inconsistent negative-displacement values.

\section{Physical units and lattice energy scale}

The model uses the lattice kinetic-energy scale
\begin{equation}
  \Ea=\frac{\hbar^2}{m_ea^2},
  \label{eq:lattice-scale}
\end{equation}
with no factor $1/2$.  For $a=1$ bohr, this convention gives one Hartree up to the
precision of the adopted physical constants.

Normalized parameters are
\begin{equation}
  t_n=\frac{T(na)}{\Ea}.
\end{equation}
They are dimensionless analysis or fitting parameters, not valid Hamiltonian
coefficients at the model boundary.  Model construction requires explicit
dimensionalization,
\begin{equation}
  T(na)=\Ea t_n,
\end{equation}
followed by any explicitly requested conversion to a physical energy unit.  The model
performs no implicit unit conversion.  Boolean values, numeric strings, unitless
Hamiltonian coefficients, and mixed coefficient units are rejected.

The model identity binds:
\begin{enumerate}[leftmargin=*]
  \item lattice constant and physical length unit;
  \item the scale convention in Eq.~\eqref{eq:lattice-scale};
  \item physical coefficient energy unit;
  \item energy-reference identity;
  \item retained basis identity;
  \item retained-space dimension; and
  \item ordered ordinary-lattice hopping inventory.
\end{enumerate}

\section{Model role and parameter acquisition}

A parameterized Bloch-hopping instance has the fixed scientific role
\texttt{REDUCED\_MODEL}.  The constructor exposes no caller-selected
control or material-reference role.  The role states that the object is an effective
parameterized model; it does not imply that any fitting or validation has succeeded.

Every parameter set must identify exactly one acquisition route:
\begin{enumerate}[leftmargin=*]
  \item \textbf{Analytic parameters}: coefficients specified by an exact mathematical
    definition;
  \item \textbf{Fourier-derived parameters}: coefficients obtained from a declared
    sampled field and transform convention;
  \item \textbf{Truncated parameters}: a declared subset of a complete ordinary-lattice
    coefficient inventory; or
  \item \textbf{Fitted parameters}: coefficients selected by a declared finite
    objective, training inventory, bounds, and optimization procedure.
\end{enumerate}

These routes are not interchangeable.  In particular, a fitted coefficient is not
reported as a native Fourier coefficient, and a support-truncated coefficient set is not
reported as a refitted model.  Bounded multistart optimization, when used, supplies a
constructive witness only; it does not establish a global optimum.

\section{Frozen scalar parent and retained field}

The predecessor scalar control derives its retained field from the four-state ambient
parent
\begin{equation}
  K(k)=T_0+e^{\ii ka}T_{+a}+e^{-\ii ka}T_{+a}^{\dagger},
  \qquad k\in\BZ,
\end{equation}
with retained scalar field
\begin{equation}
  \Hret(k)=\lambda_0\!\left(K(k)\right),
\end{equation}
where $\lambda_0$ is the least ordered eigenvalue.  The exact rational coefficients,
selected parent identity, and rejected alternatives are owned by
\path{calculations/periodic1d_rank1/band_equivalence_hopping_locality/scientific-control-decision.md}.

The ambient matrix parent has exact support $\{-a,0,+a\}$.  This fact does not imply that
the nonlinear retained eigenvalue field $\Hret(k)$ has finite ordinary-lattice support.
The parent state space, retained scalar state space, and finite coefficient
representation must remain distinct.

The retained predecessor used the exact positive reference $E_{\mathrm{ref}}=1$ and
stored dimensionless ratios to that reference.  A new execution must not silently
reinterpret those historical bytes as coefficients in electronvolts, joules, or
Hartrees.

\section{Rank-one gauge-applicability boundary}

For a scalar represented Hamiltonian and every phase field $U(k)\in U(1)$,
\begin{equation}
  U(k)^{*}H(k)U(k)=H(k).
  \label{eq:u1-triviality}
\end{equation}
Consequently, phase conjugation cannot change the represented scalar operator or its
hopping coefficients.  The calculation records one \texttt{operator-trivial}
applicability result.  It does not manufacture separate identity, global-phase,
pointwise-phase, inverse-oracle, or restricted-circuit locality channels.

Equation~\eqref{eq:u1-triviality} does not claim that a chosen eigenvector phase is
continuous, numerically stable, or physically observable.  A future study of phase
transport itself requires a separate protocol and owner.

\section{Finite cyclic and ordinary-lattice representations}

A finite discrete Fourier transform produces coefficients on a finite cyclic group.
Those coefficient classes are not automatically ordinary integer-lattice hoppings.
Every calculation must retain:
\begin{enumerate}[leftmargin=*]
  \item construction-mesh identity and ordered $k$ coordinates;
  \item forward and inverse exponent signs;
  \item transform normalization;
  \item centered cyclic representative inventory;
  \item negative-Nyquist-boundary treatment when present;
  \item explicit cyclic-to-ordinary conversion ledger; and
  \item resulting ordinary-lattice displacement inventory.
\end{enumerate}

No finite-cyclic coefficient may be silently substituted for an ordinary-lattice
coefficient.  If conversion removes, combines, or reassigns a cyclic representative,
that contribution is recorded as a separate conversion channel before interpolation or
support truncation is evaluated.

\section{Finite meshes and support inventory}

The numerical-control owner freezes the exact construction meshes, disjoint evaluation
meshes, support inventory, sampled-gap gate, application threshold, refinement rule, and
finite error--complexity diagnostic.  The current predecessor uses construction meshes
$N\in\{32,64,128,256\}$ and integer support ranges $0,\ldots,15$.  The $N=256$ label is
the largest production Fourier construction mesh, not the retained rank and not a
continuum claim.  The final mesh-stability gate compares $128\to256$.

A new protocol version is required before changing any mesh, support, threshold,
tolerance, or passage rule.  The exact frozen values are owned by
\path{calculations/periodic1d_rank1/band_equivalence_hopping_locality/numerical-control-decision.md}.

\section{Locality and support truncation}

\subsection{Descriptive locality}

For a complete ordinary-lattice coefficient inventory, descriptive locality is measured
before support removal.  For range $\Rc$, define the omitted-tail norm
\begin{equation}
  \tau_2(\Rc)
  =\left(\sum_{|R|>\Rc}|T(R)|^2\right)^{1/2}.
\end{equation}
The calculation may also retain per-displacement magnitudes and cumulative retained
power.  These values describe the selected representation; they do not themselves
construct a finite-range model.

\subsection{Deliberate finite support}

The support-truncated candidate is
\begin{equation}
  H_{\Rc}(k)=\sum_{|R|\le\Rc}T(R)e^{\ii kR}.
\end{equation}
The full declared ordinary support is retained as an excluded zero-loss control.  A
truncated candidate must identify the complete source coefficient inventory, the exact
retained support, and the exact omitted support.  Truncation is not localization, gauge
alignment, basis transformation, or refitting.

\section{Error-channel decomposition}

Let $I_N(k)$ be the complete interpolant constructed from the $N$-point mesh and let
$H_{N,\Rc}(k)$ be the ordinary-lattice candidate after deliberate support restriction.
On the disjoint evaluation mesh $E_M$, define
\begin{align}
  e_{\mathrm{int}}(N,M)
  &=\max_{k\in E_M}|\Hret(k)-I_N(k)|,\\
  e_{\mathrm{trunc}}(N,\Rc,M)
  &=\max_{k\in E_M}|I_N(k)-H_{N,\Rc}(k)|,\\
  e_{\mathrm{tot}}(N,\Rc,M)
  &=\max_{k\in E_M}|\Hret(k)-H_{N,\Rc}(k)|.
\end{align}

Interpolation error compares the direct retained field with the complete interpolant.
Truncation error isolates deliberate support deletion relative to that interpolant.
Total error compares the finite-range candidate directly with the retained field.  The
pointwise residuals telescope, but the three maxima are not added as independent
uncertainties.  Independent verification checks
\begin{equation}
  e_{\mathrm{tot}}\le e_{\mathrm{int}}+e_{\mathrm{trunc}}
\end{equation}
without replacing any stored channel by the bound.

Threshold, first-passage, and sustained-passage decisions use total error.  A direct-fit
route, if present, has its own fit residual and coefficient discrepancy; neither is
relabeled as transform interpolation or support-truncation error.

\section{Operation-first calculation sequence}

The aggregate calculation executes the following ordered workstreams:
\begin{enumerate}[leftmargin=*]
  \item \textbf{Model definition}: bind geometry, units, scale convention, basis,
    energy reference, and exact ordinary-lattice Hermiticity.
  \item \textbf{Gauge applicability}: serialize the single rank-one
    \texttt{operator-trivial} result.
  \item \textbf{Band equivalence}: construct the retained scalar field and complete
    finite cyclic Fourier representation.
  \item \textbf{Representation conversion}: produce the explicit cyclic-to-ordinary
    conversion ledger.
  \item \textbf{Hopping locality}: calculate complete-coefficient and omitted-tail
    quantities without deleting support.
  \item \textbf{Support truncation}: construct each declared finite-range candidate.
  \item \textbf{Error evaluation}: calculate interpolation, truncation-only, and total
    channels on their declared sample roles.
  \item \textbf{Aggregate decisions}: apply refinement, threshold, first-passage, and
    sustained-passage rules.
  \item \textbf{Independent verification}: reconstruct admitted values without invoking
    the production analyzer.
  \item \textbf{Authenticated rendering}: render only the closed figure inventory from
    passing digest-bound verification data.
\end{enumerate}

Reusable model storage belongs to PhysKit's immutable block-hopping model.  KSD owns the
rank-one scientific composition, configuration identity, conversion policy, thresholds,
serialization, independent-verification policy, and evidence boundary.

\section{Required immutable records}

A new aggregate execution retains canonically serialized records for:
\begin{enumerate}[leftmargin=*]
  \item protocol, configuration, calculation, and model identities;
  \item lattice constant, length unit, reciprocal convention, and sample ordering;
  \item energy-scale convention, coefficient energy unit, and energy-reference identity;
  \item parent, retained-field, and retained-basis identities;
  \item construction and evaluation meshes with explicit sample roles;
  \item complete cyclic coefficients and transform metadata;
  \item cyclic-to-ordinary conversion ledger;
  \item complete ordinary-lattice coefficients;
  \item per-range omitted-tail quantities;
  \item per-range retained and omitted support inventories;
  \item interpolation, truncation-only, total, and any separately named fit channels;
  \item finite qualification, refinement, threshold, and passage decisions;
  \item child-result identities and SHA-256 digests;
  \item independent-verification result; and
  \item authenticated visualization-data and figure digests.
\end{enumerate}

The aggregate correlates child records by exact identity and digest.  It does not copy
uncorrelated numerical values into a new result wire.

\section{Independent verification}

The verifier owns the exact admitted result bytes and must not call the production
analyzer or trust stored decisions as an oracle.  It independently reconstructs:
\begin{enumerate}[leftmargin=*]
  \item parent and retained-field samples;
  \item finite forward and inverse transforms;
  \item exact scalar Hermiticity and conjugate displacement pairing;
  \item cyclic-to-ordinary conversion;
  \item omitted-tail quantities;
  \item every declared support-truncated candidate;
  \item interpolation, truncation-only, total, and separately declared fit channels;
  \item mesh-refinement and passage decisions; and
  \item the visualization-data projection.
\end{enumerate}

Passing verification establishes agreement with this finite software protocol.  It does
not establish a theorem, continuum convergence, physical uncertainty, material
validation, or publication qualification.

\section{Lifecycle and overwrite policy}

The direct lifecycle remains:
\begin{enumerate}[leftmargin=*]
  \item \path{run_01_calculation.py}: calculate and write immutable result bytes;
  \item \path{run_02_verification.py}: independently verify the admitted result bytes;
    and
  \item \path{run_03_render.py}: render only passing authenticated data.
\end{enumerate}
Each stage refuses overwrite.  Temporary files may be used only for atomic publication.
A new operation-first aggregate uses a new execution identity and does not alter
\path{direct-controlled-four-state-n256}.

The predecessor execution remains an illustrative local unsealed finite control.  Its
wire predates the operation-first gauge-applicability, hopping-locality, and
support-truncation wrappers.  Those wrappers must be serialized in a new aggregate rather
than inserted retroactively into the predecessor result.

\section{Acceptance criteria}

A finite execution may receive the technical disposition \textbf{software-pass} only
when:
\begin{enumerate}[leftmargin=*]
  \item all required records are present, immutable, and type-valid;
  \item physical units and model identity are internally consistent;
  \item exact Hermiticity and complete transform reconstruction pass;
  \item cyclic and ordinary representations are connected by an explicit verified
    ledger;
  \item interpolation, locality, truncation, fit, and total channels remain distinct;
  \item every stored value and decision is independently reconstructed;
  \item result, verification, visualization-data, and figure digests are correlated; and
  \item direct lifecycle no-overwrite behavior passes.
\end{enumerate}

Software passage is separate from the human decisions to seal an execution, promote it
as evidence, use it in a publication, or judge a reduced model scientifically adequate.

\section{Claims and exclusions}

This protocol may support claims about the declared finite rank-one calculation,
including exact scalar $U(1)$ operator triviality, finite Fourier reconstruction,
descriptive coefficient locality, and sampled consequences of deliberate support
restriction.

It does not establish:
\begin{enumerate}[leftmargin=*]
  \item a global band-isolation theorem;
  \item continuum or infinite-mesh convergence;
  \item non-Abelian gauge behavior;
  \item a globally optimal fitted model;
  \item uncertainty quantification;
  \item equivalence to a DFT or Wannier parent;
  \item material or device validation;
  \item sealing, publication qualification, or release status; or
  \item authority to run Quantum ESPRESSO, Wannier90, remote, cluster, cloud, or other
    external computation.
\end{enumerate}

The intended output is an auditable rank-one $R$--$k$ Bloch-hopping model package in
which physical units, representation conversion, descriptive locality, deliberate
support truncation, and numerical error channels cannot be conflated.

\end{document}
